% CORRECTED VERSION
% Changes:
% - Fixed algebra in Step 7: variance term is now L η_max^2 / η_min · σ^2.
% - Updated Theorem 1 and subsequent discussion to reflect the correct bound.
% - Added an explicit justification for using η_max ≤ 1/L in Step 3.
% - Clarified the role of η_min and η_max in the bias‑variance trade‑off.
% - Marked all corrections with % CORRECTED or % ADDED comments.

\documentclass{article}
\usepackage{amsmath,amssymb,amsthm}
\usepackage{algorithm}
\usepackage{algpseudocode}
\usepackage{hyperref}
\usepackage{enumitem}
\usepackage{bm}
\newtheorem{theorem}{Theorem}
\newtheorem{lemma}{Lemma}
\newtheorem{assumption}{Assumption}
\newtheorem{remark}{Remark}
\newcommand{\EE}{\mathbb{E}}
\newcommand{\RR}{\mathbb{R}}
\newcommand{\norm}[1]{\left\lVert #1\right\rVert}
\newcommand{\inner}[1]{\left\langle #1\right\rangle}
\begin{document}

\section*{Cyclic Stochastic Gradient Descent (Cyclic SGD)}

\section*{Mathematical Formulation}
Let $f:\RR^{d}\!\to\!\RR$ be the (possibly non‑convex) objective function.  
At iteration $t\in\{0,1,\dots\}$ we draw a random sample $\xi_t$ and compute a stochastic gradient
\[
g_t \;:=\; \nabla f(w_t;\xi_t),
\]
where $w_t\in\RR^{d}$ is the current iterate.  
The learning rate follows a \emph{cyclic} schedule with period $2M$:
\[
\boxed{\;
\eta_t \;=\; \eta_{\min}
        \;+\; (\eta_{\max}-\eta_{\min})
        \Bigl(1-\Bigl|\frac{(t\bmod 2M)}{M}-1\Bigr|\Bigr),
\qquad t=0,1,\dots
\;}
\tag{1}
\]
Thus $\eta_t\in[\eta_{\min},\eta_{\max}]$ and increases linearly from $\eta_{\min}$ to $\eta_{\max}$ in the first $M$ steps, then decreases back to $\eta_{\min}$ in the next $M$ steps, repeating thereafter.

The update rule is
\[
\boxed{\, w_{t+1}= w_t - \eta_t\, g_t \, } \tag{2}
\]

\section*{Pseudocode}
\begin{algorithm}[H]
\caption{Cyclic SGD}
\begin{algorithmic}[1]
\Require Initial point $w_0$, period $M$, bounds $\eta_{\min},\eta_{\max}$, total iterations $T$
\State Set $t\gets 0$
\While{$t<T$}
    \State Compute learning rate $\eta_t$ by (1)   \Comment{triangular cyclic schedule}
    \State Sample a mini‑batch $\xi_t$ and evaluate stochastic gradient $g_t=\nabla f(w_t;\xi_t)$
    \State $w_{t+1}\gets w_t-\eta_t g_t$                \Comment{update (2)}
    \State $t\gets t+1$
\EndWhile
\State \Return $w_T$
\end{algorithmic}
\end{algorithm}

\section*{Dry Run Example}
Consider the simple quadratic $f(w)=(w-3)^2$ (so $\nabla f(w)=2(w-3)$).  
Take $w_0=0$, $\eta_{\min}=0.05$, $\eta_{\max}=0.15$, $M=2$ (period $4$).  
We perform three iterations, using the exact gradient (i.e. $g_t=\nabla f(w_t)$).

\begin{center}
\begin{tabular}{c|c|c|c}
$t$ & $\eta_t$ (from (1)) & $g_t=2(w_t-3)$ & $w_{t+1}=w_t-\eta_t g_t$ \\ \hline
0 & $\eta_{\min}=0.05$ & $2(0-3)=-6$ & $0-0.05(-6)=0.30$ \\
1 & $0.05+0.05\cdot\frac{1}{2}=0.10$ & $2(0.30-3)=-5.40$ & $0.30-0.10(-5.40)=0.84$ \\
2 & $\eta_{\max}=0.15$ & $2(0.84-3)=-4.32$ & $0.84-0.15(-4.32)=1.49$
\end{tabular}
\end{center}
Corresponding objective values:
$f(0)=9$, $f(0.30)=7.29$, $f(0.84)=4.66$, $f(1.49)=2.28$.  
The cyclic schedule already yields a faster descent than a constant $\eta=0.05$.

\newpage
\section*{Assumptions}
\begin{assumption}[Smoothness]\label{ass:smooth}
$f$ is $L$‑smooth, i.e.
\[
\forall x,y\in\RR^{d}:\qquad
f(y)\le f(x)+\inner{\nabla f(x)}{y-x}+\frac{L}{2}\norm{y-x}^{2}.
\]
\end{assumption}
\begin{assumption}[Unbiased stochastic gradients]\label{ass:unbiased}
For every $t$,
\[
\EE\!\left[\,g_t\mid w_t\right]=\nabla f(w_t).
\]
\end{assumption}
\begin{assumption}[Bounded variance]\label{ass:variance}
There exists $\sigma^{2}>0$ such that
\[
\EE\!\left[\norm{g_t-\nabla f(w_t)}^{2}\mid w_t\right]\le\sigma^{2},
\qquad\forall t .
\]
\end{assumption}
\begin{assumption}[Bounded below]\label{ass:lower}
$f$ is bounded below by $f_{\inf}>-\infty$ and we denote $ \Delta:=f(w_0)-f_{\inf}<\infty$.
\end{assumption}
\begin{assumption}[Learning‑rate bounds]\label{ass:lr}
The cyclic schedule satisfies
\[
0<\eta_{\min}\le \eta_t\le\eta_{\max}\quad\forall t,
\qquad
\eta_{\max}\le \frac{1}{L}.
\]
\end{assumption}

\section*{Main Convergence Theorem}
\begin{theorem}[Convergence of Cyclic SGD]\label{thm:convergence}
Under Assumptions \ref{ass:smooth}–\ref{ass:lr}, for any integer $T\ge 1$,
\[
\frac{1}{T}\sum_{t=0}^{T-1}\EE\!\left[\norm{\nabla f(w_t)}^{2}\right]
\;\le\;
\underbrace{\frac{2\Delta}{\eta_{\min}T}}_{\text{bias term}}
\;+\;
\underbrace{\frac{L\,\eta_{\max}^{2}}{\eta_{\min}}\;\sigma^{2}}_{\text{variance term}} .
\tag{3}
\]
Consequently,
\[
\min_{0\le t<T}\EE\!\left[\norm{\nabla f(w_t)}^{2}\right}
\;=\; O\!\Bigl(\frac{1}{\eta_{\min}T}
      +\frac{L\,\eta_{\max}^{2}}{\eta_{\min}}\sigma^{2}\Bigr).
\]
\end{theorem}
% CORRECTED: variance term now matches the algebraic derivation.

\section*{Theoretical Properties}
\begin{itemize}[leftmargin=*]
\item \textbf{Stability.} The bound (3) holds for any cyclic pattern as long as the step sizes stay within $[\eta_{\min},\eta_{\max}]$ and $\eta_{\max}\le 1/L$. Hence the algorithm is stable under the same condition as standard SGD.
\item \textbf{Hyper‑parameter sensitivity.}
    \begin{enumerate}
        \item Larger $\eta_{\max}$ \emph{increases} the variance term $\frac{L\eta_{\max}^{2}}{\eta_{\min}}\sigma^{2}$, while a larger $\eta_{\min}$ \emph{decreases} the bias term $\frac{2\Delta}{\eta_{\min}T}$.
        \item The period $M$ does not appear in (3); the bound depends only on the extremal values $\eta_{\min},\eta_{\max}$, which explains why practical performance can improve even with long cycles.
    \end{enumerate}
\item \textbf{Comparison to baselines.}
    \begin{enumerate}
        \item With a \emph{constant} step size $\eta\in(0,1/L]$, the classic SGD bound reads
        $\frac{2\Delta}{\eta T}+L\eta\sigma^{2}$.
        \item Setting $\eta_{\min}=\eta_{\max}=\eta$ in (3) recovers exactly the constant‑step bound, because $\frac{L\eta_{\max}^{2}}{\eta_{\min}}=L\eta$.
        \item For $\eta_{\max}> \eta_{\min}$ the bias term improves (since it is divided by the \emph{larger} $\eta_{\min}$) while the variance term worsens proportionally to $\eta_{\max}^{2}/\eta_{\min}$; an appropriate trade‑off can lead to a tighter bound than any fixed $\eta$.
    \end{enumerate}
\end{itemize}

\section*{Discussion}
The theorem demonstrates that, despite the non‑monotone learning‑rate schedule, the standard SGD convergence analysis carries over almost unchanged. The proof only uses the fact that each $\eta_t$ is bounded above by $1/L$ and below by a positive constant; the exact ordering of the $\eta_t$’s plays no role. Hence cyclic SGD inherits the same $O(1/T)$ decay of the average squared gradient norm (up to the additive variance term) as classical SGD, while empirically the periodic increase of the step size can help escape shallow saddles and improve practical speed.

\newpage
\section*{Preliminaries (Definitions and Lemmas)}
\begin{lemma}[Smoothness inequality]\label{lem:smooth}
If $f$ is $L$‑smooth, then for any $x$ and any stochastic gradient $g$,
\[
f(x-\eta g)\le f(x)-\eta\inner{\nabla f(x)}{g}
               +\frac{L\eta^{2}}{2}\norm{g}^{2},
\qquad\forall \eta\ge 0 .
\]
\end{lemma}
\begin{proof}
Apply the definition of $L$‑smoothness with $y=x-\eta g$ and rearrange. \qedhere
\end{proof}

\section*{Main Proof (Step‑by‑Step Derivation)}
We prove Theorem \ref{thm:convergence}. All expectations are taken with respect to all randomness up to iteration $t$ unless otherwise noted.

\noindent\textbf{Step 1. One‑step descent.}  
Using Lemma \ref{lem:smooth} with $x=w_t$, $g=g_t$ and $\eta=\eta_t$,
\begin{align}
f(w_{t+1})
&\le f(w_t)-\eta_t\inner{\nabla f(w_t)}{g_t}
      +\frac{L\eta_t^{2}}{2}\norm{g_t}^{2}. \tag{4}
\end{align}

\noindent\textbf{Step 2. Conditional expectation.}  
Take $\EE[\cdot\mid w_t]$ on both sides of (4) and invoke Assumption \ref{ass:unbiased} and the variance decomposition:
\[
\EE\!\left[\norm{g_t}^{2}\mid w_t\right]
 =\norm{\nabla f(w_t)}^{2}
   +\EE\!\left[\norm{g_t-\nabla f(w_t)}^{2}\mid w_t\right]
 \le \norm{\nabla f(w_t)}^{2}+\sigma^{2}.
\tag{5}
\]
Thus,
\begin{align}
\EE\!\left[f(w_{t+1})\mid w_t\right]
&\le f(w_t)-\eta_t\underbrace{\inner{\nabla f(w_t)}{\EE[g_t\mid w_t]}}_{\displaystyle =\norm{\nabla f(w_t)}^{2}}
   +\frac{L\eta_t^{2}}{2}\bigl(\norm{\nabla f(w_t)}^{2}+\sigma^{2}\bigr) \nonumber\\
&= f(w_t)-\Bigl(\eta_t-\frac{L\eta_t^{2}}{2}\Bigr)\norm{\nabla f(w_t)}^{2}
   +\frac{L\eta_t^{2}}{2}\sigma^{2}. \tag{6}
\end{align}

\noindent\textbf{Step 3. Lower bound on the descent coefficient.}  
Because of Assumption \ref{ass:lr} we have $\eta_t\le\eta_{\max}\le 1/L$, hence
\[
\eta_t-\frac{L\eta_t^{2}}{2}
   \;=\;\eta_t\Bigl(1-\frac{L\eta_t}{2}\Bigr)
   \;\ge\;\eta_t\Bigl(1-\frac12\Bigr)
   \;=\;\frac{\eta_t}{2}
   \;\ge\;\frac{\eta_{\min}}{2}. \tag{7}
\]
% ADDED: explicit use of $\eta_{\max}\le 1/L$ to guarantee $L\eta_t\le 1$.

\noindent\textbf{Step 4. One‑step expected decrease (with the worst‑case constants).}  
Insert (7) into (6):
\[
\EE\!\left[f(w_{t+1})\mid w_t\right]
\le f(w_t)-\frac{\eta_{\min}}{2}\norm{\nabla f(w_t)}^{2}
     +\frac{L\eta_{\max}^{2}}{2}\sigma^{2}. \tag{8}
\]
% CORRECTED: variance term correctly kept as $L\eta_{\max}^{2}\sigma^{2}/2$.

\noindent\textbf{Step 5. Removing the conditioning.}  
Take full expectation on both sides of (8):
\[
\EE\!\left[f(w_{t+1})\right]
\le \EE\!\left[f(w_t)\right]
    -\frac{\eta_{\min}}{2}\EE\!\left[\norm{\nabla f(w_t)}^{2}\right]
    +\frac{L\eta_{\max}^{2}}{2}\sigma^{2}. \tag{9}
\]

\noindent\textbf{Step 6. Telescoping sum.}  
Sum (9) for $t=0,\dots,T-1$:
\[
\EE\!\left[f(w_T)\right]
\le f(w_0)
    -\frac{\eta_{\min}}{2}\sum_{t=0}^{T-1}\EE\!\left[\norm{\nabla f(w_t)}^{2}\right]
    +\frac{L\eta_{\max}^{2}}{2}\sigma^{2}T . \tag{10}
\]
Rearrange and use $f(w_T)\ge f_{\inf}$ (Assumption \ref{ass:lower}):
\[
\frac{\eta_{\min}}{2}\sum_{t=0}^{T-1}\EE\!\left[\norm{\nabla f(w_t)}^{2}\right]
\le \Delta
   +\frac{L\eta_{\max}^{2}}{2}\sigma^{2}T .
\tag{11}
\]

\noindent\textbf{Step 7. Divide by $T$ and isolate the average.}
\[
\frac{1}{T}\sum_{t=0}^{T-1}\EE\!\left[\norm{\nabla f(w_t)}^{2}\right]
\le \frac{2\Delta}{\eta_{\min}T}
   + \frac{L\eta_{\max}^{2}}{\eta_{\min}}\sigma^{2}. \tag{12}
\]
% CORRECTED: algebra now yields the variance term $\frac{L\eta_{\max}^{2}}{\eta_{\min}}\sigma^{2}$.

Equation (12) is exactly the bound (3) claimed in Theorem \ref{thm:convergence}. \qedhere

\section*{Rate Derivation (With Constants)}
The bound (12) exhibits two terms:
\begin{itemize}
\item \emph{Bias term} $\displaystyle \frac{2\Delta}{\eta_{\min}T}$ decays as $O(1/T)$.
\item \emph{Variance term} $\displaystyle \frac{L\eta_{\max}^{2}}{\eta_{\min}}\sigma^{2}$ is independent of $T$; for a fixed target accuracy $\varepsilon$ we may choose $\eta_{\max}=O(\sqrt{\varepsilon\,\eta_{\min}})$ to make this term $O(\varepsilon)$, yielding an overall $O(1/T+\varepsilon)$ guarantee.
\item If $\sigma=0$ (full‑batch gradient) the variance term disappears and we recover the deterministic $O(1/T)$ rate for the average squared gradient norm.
\end{itemize}

\section*{Special Cases}
\begin{enumerate}[label=(\alph*)]
\item \textbf{Constant step size.} Setting $\eta_{\min}=\eta_{\max}=\eta$ reduces (12) to the classical SGD bound $\frac{2\Delta}{\eta T}+L\eta\sigma^{2}$.
\item \textbf{Purely increasing schedule.} If we restrict the cycle to the first half (monotonically increasing $\eta_t$) and stop after $M$ steps, the same proof applies with $\eta_{\max}$ replaced by the final step size; the bound still holds because only the supremum of $\eta_t$ matters.
\item \textbf{Momentum or variance‑reduction extensions.} Adding a momentum term does not affect the algebraic steps that rely solely on (1)–(2); the proof can be adapted by replacing $g_t$ with the momentum‑modified direction and bounding its variance similarly.
\end{enumerate}

\end{document}

% ===== CORRECTION SUMMARY =====
% 1. Step 7: Fixed algebraic mistake – variance term is now $\displaystyle\frac{L\eta_{\max}^{2}}{\eta_{\min}}\sigma^{2}$ instead of $L\eta_{\max}\sigma^{2}$.
% 2. Step 3: Added explicit use of $\eta_{\max}\le 1/L$ to justify $\eta_t-\frac{L\eta_t^{2}}{2}\ge\eta_t/2$.
% 3. Updated Theorem 1 and all subsequent discussion to reflect the correct bias‑variance dependence on $\eta_{\min}$ and $\eta_{\max}$.
% 4. Clarified that the bound reduces to the classic constant‑step SGD bound when $\eta_{\min}=\eta_{\max}$.
% 5. Added comments marking each correction for future reference.